% Response to reviewers -- DRAFT
% Manuscript: Cheap Features, Expensive Encoders: Rethinking Semantic Cache Decisions for LLM Systems (ESWA)
%
% Conventions
%   \reviewcomment{}  the reviewer's text, verbatim
%   \response{}       our answer
%   \change{}         what changed in the manuscript, with the location
%   \pending{}        RED: a decision or number still open; resolve before submission
% Section and table references use the labels in results/revision/manuscript_tables.tex.
% Numbers come from results/revision/COMMENT_MAP.md (five seeds unless stated).
\documentclass[11pt]{article}
\usepackage[utf8]{inputenc}
\usepackage[T1]{fontenc}
\usepackage[margin=1in]{geometry}
\usepackage{amsmath,amssymb}
\usepackage{xcolor}
\usepackage{booktabs}
\usepackage{parskip}

\definecolor{revblue}{RGB}{0,70,140}
\definecolor{chgreen}{RGB}{0,110,60}
\newcounter{reviewer}
\newcounter{rcomment}[reviewer]
\renewcommand{\thercomment}{R\arabic{reviewer}.\arabic{rcomment}}
\newcommand{\reviewer}[1]{\stepcounter{reviewer}\section*{Response to Reviewer \#\arabic{reviewer}}#1}
\newcommand{\reviewcomment}[1]{\refstepcounter{rcomment}\par\medskip
  \noindent\colorbox{revblue!8}{\parbox{\dimexpr\linewidth-2\fboxsep}{\textbf{Comment \thercomment.} \itshape #1}}\par}
\newcommand{\response}[1]{\par\noindent\textbf{Response.} #1\par}
\newcommand{\change}[1]{\par\noindent{\color{chgreen}\textbf{Change in the manuscript.} #1}\par}
\newcommand{\pending}[1]{{\color{red}\textbf{[#1]}}}
% Until the manuscript numbers are final, references show the manuscript label.
% For the final letter, load the manuscript's .aux with the xr package and delete this line.
\renewcommand{\ref}[1]{\textsf{\small[\detokenize{#1}]}}

\begin{document}

\begin{center}
{\Large\bfseries Response to the Reviewers}\\[4pt]
Cheap Features, Expensive Encoders: Rethinking Semantic Cache Decisions for LLM Systems\\
Expert Systems with Applications \pending{manuscript number}
\end{center}

Dear Editor and Reviewers,

We thank the reviewers for careful and constructive reports. Their comments led us to re-run the whole
evaluation and, in two places, to change what the paper claims. Below, each comment is quoted in full and
followed by our response and the corresponding change in the manuscript.

\subsection*{Summary of the main changes}
\begin{enumerate}
  \item \textbf{Corrected stream labels.} While re-running the experiments we found that the stream generator
    could give the same distinct question different labels in different requests, and that distractors could
    recur verbatim. Both are fixed: every distinct question now has one stable label, and distractors are drawn
    from strings never emitted before. On the same cache decisions, the corrected labels raise hit precision
    from 75--77\% to 89--94\%; the hit decisions themselves are unchanged. All stream results are re-reported
    with the corrected generator, and the correction is stated in the evaluation section.
  \item \textbf{Five independent streams and statistical tests.} Every stateful result is now reported as a mean
    with a 95\% confidence interval over five independent streams (seeds 42--46); a sixth stream (seed 1042) is
    reserved for tuning. Comparisons use seed-level paired tests with Holm correction.
  \item \textbf{Real generation waits on GPU hardware.} End-to-end timing now replays answers from
    \texttt{gpt-5.4-nano} with their real generation waits, on an NVIDIA RTX 4090, and adds latency distributions,
    top-$k$ retrieval, cross-encoder batching and concurrency.
  \item \textbf{False hits are priced.} We add false-discovery rate, a risk-adjusted saving
    $\mathrm{RAS}(\lambda)$, the break-even false-hit cost, full threshold sweeps for C1--C4, abstention, and a
    selectively gated cross-encoder (new configuration C5).
  \item \textbf{A clearer central claim.} At equal served-hit ratio, C1--C4 lie on one precision--coverage
    curve; choosing the threshold for the cost of a false hit matters far more than which configuration is used.
    The cross-encoder raises cost without moving that curve. C2 is retained as the default because it offers the
    most precise operating point at its default threshold for 1.4$\times$ the cost of C1, not because it is a
    better classifier. The recommendation is now stated in these terms.
  \item \textbf{Semantic Decay is reported as an exploratory policy with a negative result}, with capacities
    200--1000, five streams, tuned and untuned weights, a two-term variant, SphereLFU (2026) as a further
    baseline, eviction-time density statistics, and a control showing that its precision gain is matched by the
    decision threshold.
  \item \textbf{New baselines and data.} vCache (ICLR 2026) is added as a baseline and SphereLFU (2026) as an
    eviction baseline; a real search-query log (SemBenchmarkSearchQueries, 50,000 queries) and a
    10,000-request stream are added; an LLM judge audits answer substitutability.
\end{enumerate}

%=====================================================================================================
\reviewer{}

\reviewcomment{Why C2 rather than C1 as the default? [\ldots] So on the very metrics the paper cares
about---latency and correct token saving---C1 looks like the better cheap default; C2 only wins on hit precision
by 0.8 points and on hard-negative rejection. The authors need to be explicit about why C2 is preferred. Is it
mainly because C2 rejects near-duplicate distractors better? If so, that should be stated as the criterion, not
left implicit. Right now the recommendation reads as if it is driven by the stacking framework's novelty rather
than by the serving objective.}

\response{We agree that the criterion was left implicit, and the re-run evaluation changes the answer in an
instructive way. With corrected labels, five streams and real generation waits
(Table~\ref{tab:end_to_end}, Table~\ref{tab:seed_variation}), C1 serves 2.25 points more hits than C2 and
0.63 points more correct-only token saving, and costs 2.9~ms less per request; C2 is 2.7 points more precise
(94.9\% vs.\ 92.2\%), so its false-discovery rate is roughly a third lower. All of these differences are
significant (Table~\ref{tab:c1_c2_paired}). The reviewer's reading of the trade-off is therefore correct: C1
wins on coverage and cost, C2 on precision.

The decisive evidence comes from sweeping each configuration's threshold (Table~\ref{tab:matched_precision}):
at the same served-hit ratio, C1, C2, C3 and C4 are statistically indistinguishable in precision, and none of
them beats C1 once each threshold is chosen for the false-hit cost $\lambda$ (Table~\ref{tab:best_per_lambda};
no difference survives Holm correction). What differs is where each configuration sits at its default
threshold. At $\lambda=10$, choosing C1's threshold for the cost raises its risk-adjusted saving from 6.3\% to
41.4\%, while the differences between configurations are about one point.

The criterion is therefore the serving objective once false hits have a price. C2's default threshold already
sits at a precise operating point (break-even false-hit cost 18.9 vs.\ 11.9 for C1), which makes it the safer
default when the cost of a wrong answer is unknown. An operator who knows that cost can calibrate C1's threshold
and obtain the same result for less compute. Hard-negative rejection is part of this, since most false hits in
our streams are near-duplicate distractors, but the general statement is about the operating point. We have
rewritten the recommendation accordingly, and it no longer rests on the stacking framework.}

\change{Abstract, Introduction (contributions), Sec.~Stateful Semantic Caching Performance (new
Tables~\ref{tab:seed_variation} and \ref{tab:c1_c2_paired}), new subsection ``Pricing False Hits''
(Tables~\ref{tab:risk_adjusted}, \ref{tab:matched_precision}, \ref{tab:best_per_lambda}), Discussion and
Conclusion: the default is justified by precision at the default operating point and by the cost of false
hits; the paper states that a calibrated C1 threshold is equivalent when $\lambda$ is known.}

\reviewcomment{False hits are identified but not priced. [\ldots] I would like to see false discovery rate or
false-hit rate reported alongside hit precision, and some risk-adjusted notion of saving, for example effective
saving = correct-hit saved tokens $-\ \lambda\ \times$ false-hit cost. A threshold sweep would also help: how do
precision, correct-only saving, and latency move as the LightGBM or bi-encoder threshold changes? [\ldots] A
short discussion of abstention, secondary verification, or using the cross-encoder only on high-risk candidates
would also strengthen the practical relevance.}

\response{We have adopted the reviewer's proposal directly. First, the 75--77\% precision in the submitted
version was largely a labelling artefact (summary item~1); with corrected labels, precision is 92--95\%, so
about one served hit in 13 to 21 is wrong rather than one in four. That is still too high to leave unpriced.

We now report the false-discovery rate and false-hit rate for every configuration, and define the
risk-adjusted saving $\mathrm{RAS}(\lambda) = (T_{\text{correct}} - \lambda\, T_{\text{false}})/T_{\text{total}}$,
where each false hit costs $\lambda$ times the tokens of the call it replaced, together with the break-even
$\lambda$ at which caching stops paying (Table~\ref{tab:risk_adjusted}). At $\lambda=10$, the default C1 keeps
9.8\% of the tokens and C2 25.0\%. Full threshold sweeps for C1--C4 (bi-encoder cosine, LightGBM probability
and cross-encoder score) give precision, correct-only saving and $\mathrm{RAS}$ along each curve
(Tables~\ref{tab:matched_precision} and \ref{tab:best_per_lambda}). The decision cost does not depend on the
threshold, because the same computation runs for every request; response time changes only through the hit
ratio.

On the practical side, we evaluate (i) abstention, which sends a request whose meta-classifier probability is
within a margin of the threshold to the LLM (precision rises from 94.8\% to 96.2\% for a 1.6-point loss in hits),
and (ii) a new configuration C5 that calls the cross-encoder only on uncertain candidates, i.e.\ ``using the
cross-encoder only on high-risk candidates'' (Table~\ref{tab:gating}). C5 calls the cross-encoder on 1--14\% of
requests and costs 9.8--11.2~ms instead of 21.7~ms. On our streams it adds up to 1.4 points of hits for a
0.1--1.8-point loss in precision, so it does not improve on C2. Selective verification is therefore cheap to
deploy, but on this workload it does not change the conclusion.}

\change{Sec.~Evaluation Metrics (definitions of FDR, $\mathrm{RAS}(\lambda)$ and break-even $\lambda$); new
subsection ``Pricing False Hits''; new subsection ``Top-k Retrieval and Selective Verification''
(Tables~\ref{tab:topk}, \ref{tab:gating}); C5 added to Sec.~Evaluated Cache Configurations.}

\reviewcomment{The 60 ms cross-encoder cost is not general enough. [\ldots] At minimum, the paper should
discuss cross-encoder cost under top-k re-ranking [\ldots] and under GPU batching where the per-pair cost may
drop. Latency distributions under concurrency would also be more informative than the mean alone. [\ldots] the
online recommendation should be scoped explicitly to the hardware and batch-size-1 setting used here.}

\response{We re-measured every latency on an NVIDIA RTX 4090 and added the analyses the reviewer lists.
\begin{itemize}
  \item \emph{Per-stage cost} (Table~\ref{tab:latency_breakdown}): at batch size 1 the cross-encoder takes
    11.5~ms per pair, 61\% of C3's decision path; bi-encoding takes 6.9~ms and the cheap features plus LightGBM
    2.9~ms.
  \item \emph{Batching} (Table~\ref{tab:cross_batching}): the per-pair cost falls from 11.0~ms at batch 1 to
    1.3~ms at batch 16. Batching helps re-ranking several candidates of one request, but not a single online
    request, which needs one pair scored immediately.
  \item \emph{Top-$k$} (Table~\ref{tab:topk}): scoring $k=10$ candidates raises C3's cost from 18.2 to 21.0~ms;
    C2 stays at 10~ms.
  \item \emph{Distributions and concurrency} (Table~\ref{tab:latency_distribution}): p50/p95/p99 at batch~1,
    and 1, 4 and 16 concurrent clients. At 16 clients, C4's p99 is 512~ms against 269~ms for C2, and its
    throughput falls to 38 requests/s against 87.
\end{itemize}
The cross-encoder's absolute cost is much lower on the GPU than the 60~ms measured on the M2, but its share of
the decision path, and its effect under concurrency, remain. On a hit it lies on the critical path, and it can be
overlapped with generation only by starting the LLM call speculatively on every request. We scope the recommendation explicitly: on the hardware we
measured, at batch size 1 and with one request in flight, the cross-encoder does not pay for itself in the
online path; with large batches or idle accelerator capacity its marginal cost shrinks, and gating (C5) is then
the natural way to use it.}

\change{Sec.~Latency Breakdown rewritten with GPU measurements (Tables~\ref{tab:latency_breakdown},
\ref{tab:latency_distribution}, \ref{tab:cross_batching}); Sec.~Experimental Setup reports the hardware;
Discussion and Conclusion scope the recommendation to the measured setting. \pending{CHECK: the end-to-end
speed-up is now 2.1--2.2$\times$ instead of about 4$\times$; state the reason (real generation waits of
\texttt{gpt-5.4-nano}) in Sec.~Stateful Semantic Caching Performance.}}

\reviewcomment{Semantic Decay is exploratory, not a validated gain. [\ldots] I would not present Semantic Decay
as a contribution with confirmed value. It should be framed as an exploratory policy. The authors should report
results over more capacities, independent streams, and seeds, and use statistical tests against LRU/LFU rather
than single-point comparisons. Given the age term adds little, a two-term version---freshness plus boost---may be
simpler and equally effective. The fact that the weights were chosen on the same workload family also raises a
tuning-leakage concern that should be acknowledged.}

\response{We agree, and the extended evaluation confirms the reviewer's reading. Semantic Decay is now presented
as an exploratory policy, and the paper reports a negative result for it (Table~\ref{tab:eviction_capacity}).
\begin{itemize}
  \item \emph{Design.} Nine capacities (200--1000), five independent evaluation streams, seed-level paired tests
    against LRU, LFU and SphereLFU (a 2026 semantic eviction policy) with Holm correction.
  \item \emph{Tuning leakage.} The weights are now tuned on a separate stream (seed 1042) and frozen before
    evaluation. We also state that the weights in the submitted version (0.28/0.38/0.34) were set by hand and
    were not the result of a search.
  \item \emph{Two-term variant.} The tuned three-term weights are $(0,0,1)$: only the boost term survives, and
    the freshness-plus-boost variant tunes to the same policy, confirming that the age and freshness terms add
    nothing.
  \item \emph{Result.} After Holm correction, Semantic Decay is not significantly better than LRU at any
    capacity and is worse than LFU at capacities 200--300. Frequency, not semantic similarity, is what helps.
  \item \emph{Precision.} At capacities 200--400 the tuned policy serves 16--40\% fewer false hits than LFU. A
    control experiment shows that this is not a gain: once each policy's decision threshold is also chosen for
    $\lambda$ on the tuning stream, no difference from LFU or SphereLFU survives Holm correction. The policy
    serves fewer, safer hits, and the threshold buys the same trade.
\end{itemize}
The control also explains where the false hits come from: hits on entries far from the popular region of the
cache are false 14.7\% of the time, against 1.0\% for the most central entries. We report this as a
diagnostic finding.}

\change{Introduction (Semantic Decay removed from the list of contributions and described as exploratory);
Sec.~Eviction Policy Evaluation rewritten with Table~\ref{tab:eviction_capacity}, the tuning protocol, the
threshold control, and the provenance of the original weights; limitation on tuning leakage added.
\pending{DECIDE: whether to report the admission-control result (frequency-gated admission adapted from
TinyLFU, +1.4 to +3.7 correct-hit points over LFU at capacities 200--600) in one sentence or omit it.}}

\reviewcomment{The workload is still synthetic. [\ldots] I would not require a full production trace for a major
revision, but at least a semi-real query log, a multi-turn dialogue scenario, or a sensitivity analysis over Zipf
exponent, repetition rate, and hard-negative ratio would make the conclusions more credible. Privacy and
cross-user isolation are also worth a brief mention.}

\response{We added three of the reviewer's suggestions.
\begin{itemize}
  \item \emph{Sensitivity analysis} (Table~\ref{tab:workload_sensitivity}). We varied the Zipf exponent
    (0.8--1.4), repetition rate (0.3--0.7) and hard-negative rate (0.1--0.6), one factor at a time, over three
    streams per setting. The ordering of the configurations does not change: C2 and C4 are more precise than C1
    in every cell. The repetition rate dominates both hit ratio and precision.
  \item \emph{Semi-real query log} (Tables~\ref{tab:searchqueries} and \ref{tab:searchqueries_matched}). We use
    the first 50,000 queries of SemBenchmarkSearchQueries, real Bing search queries (ORCAS) grouped into
    equivalence classes by the vCache authors. With thresholds and models tuned on QQP, every configuration loses
    most of its precision (C1 41\%, C2 53\%): QQP-tuned decision boundaries do not transfer to search traffic.
    After retraining on queries from groups that never appear in the test stream, C2 is never better than a
    re-tuned C1 at the same hit ratio (between $-1.8$ and $+0.1$ points). This strengthens our central
    conclusion: the operating point must be calibrated per domain, and cheap features do not replace that
    calibration.
  \item \emph{Long stream.} A 10,000-request stream (Table~\ref{tab:long_stream}) shows stable hit ratios across
    1,000-request blocks.
\end{itemize}
We did not add a multi-turn dialogue workload, because we have no public multi-turn trace with equivalence
labels; the ContextCache comparison therefore remains an evaluation of its decision rule on synthesized
sessions, and we state this limitation explicitly. We added a paragraph on privacy and cross-user isolation:
serving one user's cached answer to another can leak personal content, so a deployed cache should partition
entries per tenant or restrict sharing to answers that contain no user data, at some cost in hit ratio.}

\change{New subsubsection ``Workload Sensitivity''; new subsection ``Real Query Log'' in Sec.~Cross-Dataset
Validation (and the dataset in Sec.~Datasets); long-stream results in Sec.~Comparison with Published Semantic
Caches; Limitations: multi-turn traffic, privacy and cross-user isolation.}

%=====================================================================================================
\reviewer{}

\reviewcomment{The organization and logical structure of the Related Work section need to be further improved.
[\ldots] The authors are encouraged to provide a more systematic classification and synthesis of the existing
studies, and to further clarify the limitations of prior work and the distinctions between the proposed method
and existing approaches.}

\response{We reorganised Related Work around the three decisions a semantic cache makes, and each part now
ends with the gap this paper addresses:
(i) \emph{whether to serve a hit}: fixed similarity thresholds (GPTCache, SCALM), learned or per-entry thresholds
(VectorQ, vCache, MVR-cache) and LLM-based verification (LaCache);
(ii) \emph{what context to match}: conversation-aware caching (ContextCache) and tiered designs (Krites);
(iii) \emph{what to keep}: eviction and admission (LRU/LFU, TinyLFU, SphereLFU, RAC, recent LLM-cache eviction
studies, cost-learned content selection, and SOLAR).
A summary table compares the methods by decision signal, cost per request, need for labels or feedback, and
whether they handle capacity. The distinction of our work is stated in one place: we measure what each decision
signal buys \emph{per unit of serving latency} in a stateful cache, and price false hits explicitly.}

\change{Sec.~Related Work restructured with a comparison table. \pending{TO WRITE}}

\reviewcomment{There is an inconsistency between the description and the illustration in Figure 1. The caption
states that ``The cross-encoder (dashed) is active only in C3 and C4.'' However, the corresponding dashed
Cross-encoder module and the distinction between C3 and C4 do not appear to be clearly shown in Figure 1.}

\response{Thank you for catching this. Figure~1 has been redrawn: the cross-encoder is drawn with a dashed
outline, and a small activation strip under the diagram marks which components each configuration (C1--C5)
uses, so C3 (bi-encoder + cross-encoder) and C4 (all components) are visibly distinct. The caption now matches
the figure.}

\change{Figure~1 and its caption. \pending{TO DRAW}}

\reviewcomment{The overhead of the Full mode in Table 9 requires further explanation. [\ldots] provide a more
detailed analysis of the computational cost associated with each component [\ldots] and to discuss the
trade-off between the additional computational cost and the corresponding performance improvement.}

\response{Table~\ref{tab:latency_breakdown} now breaks the decision path into bi-encoding, vector search,
cross-encoding and features plus LightGBM, measured on the GPU at batch size~1. The Full mode (C4) costs
21.6~ms: 6.9~ms bi-encoding, 0.3~ms search, 11.5~ms cross-encoding and 2.9~ms features plus LightGBM. The
cross-encoder alone accounts for 53\% of it, and is the reason C4 costs 3.0$\times$ C1 and 2.1$\times$ C2.
On the benefit side, C4 is the most precise configuration at its default threshold (95.2\% vs.\ 94.9\% for C2), but at equal served-hit ratio it is indistinguishable from C1 and C2
(Table~\ref{tab:matched_precision}). The extra 11.5~ms therefore buys no improvement that a cheaper
configuration with a different threshold cannot match. If a cross-encoder is used, the gated configuration C5
keeps most of C2's cost (10.7~ms) by calling it on a small fraction of requests.}

\change{Sec.~Latency Breakdown (Table~\ref{tab:latency_breakdown} and the paragraph after it); the cost
columns of the ablation tables updated to the GPU measurements.}

\reviewcomment{The experimental comparison with recent methods published in 2026 should be strengthened. [\ldots]
include representative state-of-the-art methods published in 2026 as additional baselines and conduct fair
comparisons under consistent experimental settings.}

\response{We added vCache (Schroeder et al., ICLR 2026), which learns a per-entry threshold with an error bound
$\delta$, as a decision-layer baseline, following the authors' implementation. We also added SphereLFU (Biton
and Friedman, 2026) as an eviction baseline. Both run on the same streams, encoder, answer pool and hardware as
our configurations (Tables~\ref{tab:baseline_comparison}, \ref{tab:baseline_timing}, \ref{tab:long_stream},
\ref{tab:searchqueries}, \ref{tab:eviction_capacity}).

vCache is the most precise system we evaluated (98.5--99.5\%), but serves only 10--28\% hits against
51--54\% for the threshold-based systems. It also needs a correctness signal for most requests (0.72--0.90
judge calls per request in our runs, supplied here by the oracle labels), which a deployment would have to pay
for. On the real query log it serves 1--4\% hits. SphereLFU improves on LRU at small capacities but not on LFU.

Other 2026 methods are discussed in Related Work with the reason they were not run: MVR-cache reduces to vCache
on single-sentence queries; LaCache requires an LLM decoding step per candidate; Krites requires a curated
static tier; RAC \pending{DECIDE: cite only, or add as an eviction baseline}.}

\change{Sec.~Comparison with Published Semantic Caches (vCache added to the systems and tables); SphereLFU in
Sec.~Eviction Policy Evaluation; Related Work.}

\reviewcomment{The manuscript requires further English language polishing.}

\response{The manuscript has been revised throughout for grammar, wording and transitions. Long sentences have
been split, terminology unified (e.g.\ ``served hit'', ``false hit'', ``operating point''), and every section
now opens with a sentence that states its purpose.}

\change{Throughout. \pending{TO DO LAST}}

%=====================================================================================================
\reviewer{}

\reviewcomment{The cache decision uses only the top-1 retrieved candidate; many production pipelines consider
top-k candidates to mitigate single-retrieval error and to enable selective cross-encoder gating on uncertain
cases.}

\response{The pipeline now retrieves the top-$k$ candidates, scores every one with the configuration's decision
rule, and serves the best accepted candidate (Table~\ref{tab:topk}, $k\in\{1,3,5,10\}$). On our streams, $k>1$
changes hit ratio by at most 0.1 points and lowers precision by 0.3--0.8 points for the configurations that
re-rank. The extra candidates are almost always non-equivalent, and scoring them gives more chances to accept a
wrong one. The cost grows with $k$ for the cross-encoder configurations (C3: 18.2 to 21.0~ms) but hardly at all
for C2. Selective gating is implemented as the new configuration C5 (Table~\ref{tab:gating}): the
meta-classifier decides alone outside an uncertainty band, and the cross-encoder decides inside it. The band is
chosen on out-of-fold probabilities, and an empty band reproduces C2 exactly.}

\change{Sec.~Query Processing (top-$k$ scoring); C5 in Sec.~Evaluated Cache Configurations; new subsection
``Top-k Retrieval and Selective Verification''.}

\reviewcomment{Semantic Decay yields limited improvements and only at a single capacity point; policy is
fixed-parameter and not adaptive to workload shift, contrasting with recent learning-augmented
eviction/admission designs.}

\response{We agree; see also our response to Comment~R1.4. Across nine capacities and five streams, Semantic
Decay is not significantly better than LRU and is worse than LFU at small capacities, and the paper now says so.
Its parameters are tuned on a separate stream and then fixed. We now state explicitly that the policy is not
adaptive, and we compare it with SphereLFU, a 2026 semantic eviction policy with online, decaying frequency
estimates. \pending{DECIDE: add LeCaR, a learning-augmented eviction baseline (implemented; run
pending), and a workload-shift experiment, or discuss adaptivity as future work only.} The learning-augmented
and admission designs the reviewer refers to are discussed in the restructured Related Work.}

\change{Sec.~Eviction Policy Evaluation; Related Work (capacity management); Limitations.}

\reviewcomment{Baseline reimplementations (GPTCache, SCALM, VectorQ, ContextCache) are mentioned but
quantitative comparisons are not shown in the provided text; mapping conversation- or pattern-aware methods onto
QQP may be unfair without careful protocol alignment.}

\response{Table~\ref{tab:baseline_comparison} now gives the quantitative comparison over five streams: hit
ratio and precision with confidence intervals, correct-only saving, break-even $\lambda$, and a paired test
against C2. Table~\ref{tab:baseline_timing} gives their timing with real generation waits.

On protocol alignment, every baseline is a re-implementation of the published \emph{decision rule} on our shared
encoder, index and answer pool. The comparison isolates the decision, not the surrounding system, and each rule
is calibrated on the same validation split as C1. We now state this in the text. For ContextCache, whose
conversation-aware rule needs dialogue context, QQP has no natural sessions; we evaluate the rule on synthesized
sessions and say plainly that this tests the rule, not the setting it was designed for. Threshold sweeps
(Figure~F3) show that GPTCache, SCALM, VectorQ and ContextCache lie on the same precision--coverage curve as C1,
so their differences from C2 at default thresholds come from the operating point.}

\change{Sec.~Comparison with Published Semantic Caches (protocol paragraph, Tables~\ref{tab:baseline_comparison}
and \ref{tab:baseline_timing}, operating-curve figure).}

\reviewcomment{Could you report cache-density or redundancy statistics (e.g., neighbor counts within the hit
radius) at eviction time to contextualize why Semantic Decay helps (or not) at certain capacities?}

\response{Yes. Table~\ref{tab:eviction_density} reports, at each eviction, the number of live entries within the
hit radius of the victim, the share of victims whose topic is still covered by another live entry, and the share
of victims that were never hit. Redundancy is low in every policy: at capacity 400 a victim has on average
0.035--0.056 neighbours within the hit radius, and its topic is still covered in only 1.4--2.2\% of evictions. Semantic Decay
is built to evict redundant entries, but the cache contains very few of them, which explains why it cannot
outperform frequency-based policies. The diagnostic in our response to Comment~R1.4 adds the complementary
picture: the entries that matter for precision are the isolated ones, and hits on them are false about 15\% of
the time.}

\change{Sec.~Eviction Policy Evaluation, new Table~\ref{tab:eviction_density} and the paragraph interpreting it.}

\reviewcomment{could with additional experiments, especially quality-adjusted hit audits, selective
cross-encoder gating, top-k retrieval, encoder/ANN sensitivity, and stronger baseline comparisons.}

\response{All five are now in the paper:
quality-adjusted hit audit (Table~\ref{tab:answer_audit}, see Comment~R3.7);
selective gating (C5, Table~\ref{tab:gating});
top-$k$ retrieval (Table~\ref{tab:topk});
encoder and ANN sensitivity (Tables~\ref{tab:encoders} and \ref{tab:index}): three bi-encoders
(all-mpnet-base-v2, all-distilroberta-v1, all-MiniLM-L6-v2), each with its own threshold, preserve the
precision advantage of C2 over C1 (+2.3 to +3.4 points), and HNSW indexing with efSearch 8--64 gives the same hit
ratio and precision as exact search;
stronger baselines (vCache and SphereLFU, see Comment~R2.4).}

\change{New subsections ``Answer-Level Audit'' and ``Encoder and Index Sensitivity''; see the changes listed for
Comments R2.4, R3.1 and R3.7.}

\reviewcomment{The workload is derived from QQP paraphrase pairs; answer substitutability can differ materially
from paraphrase similarity. While ``hit precision'' is reported, it is unclear how answer
correctness/substitutability is adjudicated beyond QQP labels.}

\response{We now state how correctness is adjudicated (a hit is correct when the served entry and the request
belong to the same QQP equivalence class), and we audit it at the answer level. For every served hit of C1, C2
and C4 on one stream, an LLM judge (\texttt{gpt-6-luna}) is given the request, the cached answer and a freshly
generated reference answer, and decides whether the cached answer is an acceptable substitute
(Table~\ref{tab:answer_audit}). Judge
precision is within 1.5 points of label precision for every configuration (C1 91.8\% vs.\ 92.1\%; C2 96.1\%
vs.\ 94.8\%; C4 93.6\% vs.\ 95.0\%), and C2 has the highest judged precision. Agreement between judge and labels
is moderate (Cohen's $\kappa=0.40$): the judge accepts 48 label-false hits as substitutable and rejects 14
label-correct ones. Label precision is therefore a reasonable, slightly conservative proxy for answer
substitutability on this workload. We report the judge prompt, the shared 256-token answer cap, and the
calibration on 30 pairs.}

\change{Sec.~Evaluation Framework (how correctness is adjudicated; judge model and prompt; answer cap); new
subsection ``Answer-Level Audit'' (Table~\ref{tab:answer_audit}).}

\bigskip
We thank the reviewers again; their comments have made the paper's claims narrower but much better supported.

\end{document}
